Frozen graph classifiers expose a practical question for low-timestep conversion: does matching final graph logits improve calibration beyond reconstructing local activations? We compare six conditions using the same pretrained residual GCNs, training examples and update budget on MNIST and CIFAR10 superpixel classification. Matched membrane conditions calibrate 693 initial-membrane values while retaining hard signed eight-level events. Across three teacher seeds per task, frozen test differences between graph-logit and local reconstruction objectives range from $-0.18$ to $+0.20$ percentage points at one or two active ticks; validation and test results differ in direction. These official splits have historical project exposure, so the endpoints are descriptive. Output error correlates with prediction flips more closely than node error, while ideal-site interventions locate a substantial short-prefix deficit upstream: restoring the first node nonlinearity alone closes about 54\% and 24\% of the one-tick teacher-disagreement gaps. Graph-level objective choice consequently requires evidence beyond error correlation. Synchronized RTX 3060 measurements show that the current diagnostic conversion core takes about 2.4--3.1 times the teacher's time at one tick. The study separates fidelity, calibration effectiveness and software execution cost, with small objective-dependent gains and no general method advantage established in the tested regime.
